\documentclass{article}
\usepackage{amsmath}
\begin{document}

\textbf{Summary of Strassen's Algorithm and Implementation Notes}

\section*{Conceptual Overview}
1. \textbf{Recursive Partitioning}:
    - Strassen’s algorithm recursively divides matrices \( A \) and \( B \) into four quadrants: 
      \[
      A = \begin{bmatrix} A_{11} & A_{12} \\ A_{21} & A_{22} \end{bmatrix} \quad B = \begin{bmatrix} B_{11} & B_{12} \\ B_{21} & B_{22} \end{bmatrix}
      \]
    - This partitioning continues until reaching the base case of \( N = 1 \), where scalar multiplication is performed.

2. \textbf{Base Case}:
    - When \( N = 1 \), each quadrant is a single element of the original matrix, and the actual multiplication occurs:
      \[
      C[0][0] = A[0][0] \times B[0][0]
      \]
    - The recursion then unwinds, merging results from smaller matrices back up to form the final matrix \( C \).

3. \textbf{Combining Results}:
    - The results of the multiplications are combined as the recursion unwinds:
    \[
    \begin{aligned}
    &C_{11} = M_1 + M_4 - M_5 + M_7 \\
    &C_{12} = M_3 + M_5 \\
    &C_{21} = M_2 + M_4 \\
    &C_{22} = M_1 - M_2 + M_3 + M_6 
    \end{aligned}
    \]
    - This recombination reduces the number of multiplications from \( N^3 \) to \( N^{\log_2 7} \approx N^{2.81} \).

4. \textbf{No Explicit Partitioning in Code}:
    - The algorithm conceptually partitions matrices, but in practice, no new matrices are created.
    - Pointers, indices, or views are used to reference submatrices, avoiding overhead from copying.

\section*{Implementation and Optimization Notes}

1. \textbf{Zero-Based Indexing in Code}:
    - Instead of using \( A_{11}, A_{12} \), etc., the submatrices can be accessed using zero-based indexing such as \( A[0][0], A[0][1], \ldots \).

2. \textbf{Custom Subscript Operator}:
    - Overloading the subscript operator can be useful for managing matrix partitioning. However, it should be evaluated for its trade-offs, as it might conflict with data-oriented design principles.

3. \textbf{Parallelization and SIMD}:
    - Given the recursive structure and power-of-2 dimensions, Strassen’s algorithm has potential for optimization with SIMD instructions or multi-threading.
    - Consider using row-major order for better cache efficiency and easier SIMD vectorization.

4. \textbf{Base Case Optimization}:
    - The base case of \( N = 1 \) can be a critical bottleneck. Using intrinsic functions or hardware acceleration at this level could significantly improve performance.

5. \textbf{Understanding the Data Flow}:
    - Before implementing optimizations, it’s essential to have a clear understanding of how the algorithm flows and how results are collated as the recursion unwinds.
    - The recurrence relations and the way intermediate products are added back into \( C \) are key to implementing the algorithm correctly.

\end{document}
